% ECONOMICS_PROBLEM: {"id":"C73-17","title":"Maxmin equality with the concavified nonrevealing value","jel_code":"C73","source_id":"OP-0168"}
% FIRST_POSTED: 2026-10-10
% ATTEMPT_STATUS: solved
% ENTRY_KIND: research_attempt
% Proof of the full bounded-Borel, all-prior target.
% Established inputs: Blackwell Borel determinacy and Borel superhedging duality.
\documentclass[11pt,a4paper]{article}
\usepackage[margin=20mm]{geometry}
\usepackage[T1]{fontenc}
\usepackage{lmodern,microtype}
\usepackage{amsmath,amssymb,amsthm,mathtools}
\newtheorem{theorem}{Theorem}
\usepackage{xurl,hyperref}
\hypersetup{hidelinks,hypertexnames=false,
pdftitle={C73-17: An exact local-geometry characterization of maxmin concavification}}
\setlength{\parindent}{0pt}
\setlength{\parskip}{3pt}
\setlength{\emergencystretch}{2em}
\allowdisplaybreaks
\raggedbottom
\widowpenalty=10000
\clubpenalty=10000
\displaywidowpenalty=10000
\title{C73-17: An exact local-geometry characterization\\of maxmin concavification}
\author{Ji Ho Bae}
\date{10 October 2026}
\begin{document}
\maketitle

\section*{Model, process, and claim status}
\textbf{Model:} GPT Codex. \textbf{Contributor:} Ji Ho Bae.
The exact serving revision was not exposed. The mathematical argument was
produced with parallel derivation and adversarial checking in response to a
request for a complete hand proof of the exact catalogue statement.
No human-written mathematical edits are included. This is an LLM claim of a
complete necessary-and-sufficient characterization for the full bounded-Borel,
all-prior target; community expert review is pending. No proof-assistant
verification is claimed.
The complete original generated output is preserved at the immutable
\href{https://github.com/jbaelaw/maxmin-concavification-characterization/blob/448392472b35ef2303148f78fbbd819147bab774/paper/Proof.tex}{manuscript TeX}
and \href{https://github.com/jbaelaw/maxmin-concavification-characterization/blob/448392472b35ef2303148f78fbbd819147bab774/paper/Proof.pdf}{manuscript PDF}.

The criterion is an infinitary local-capital certificate. Conserved likelihood
flows encode the informed player's likelihoods; explicit behavioral action
kernels, expected-payoff optimization, and game values are absent from its
definition. Blackwell Borel determinacy is used only for nonrevealing games,
and Borel superhedging duality for all-path capital bounds. No value of the
incomplete-information game is asserted.

\section{Original target and notation}

Fix nonempty finite sets \(K,I,J\), put \(m=|K|\), \(X=I\times J\), \(H=\bigcup_{t\ge0}X^t\), and \(\Omega=X^{\mathbb N}\). The empty history is \(\varnothing\), \(hij\) means append the pair \((i,j)\), and \(\omega|t\) is the length-\(t\) prefix. Let \(f:K\times\Omega\to\mathbb R\) be bounded and Borel. The informed maximizer has behavioral prescriptions \(\sigma_k(i\mid h)\); the uninformed minimizer has prescriptions \(\tau_j(h)\), with simultaneous independent current randomizations and full observation of past joint actions.

Write \(L_f(p)=\sup_\sigma\inf_\tau\mathbb E_{p,\sigma,\tau} f\). The nonrevealing payoff is \(f_q(\omega)=\sum_kq_k f(k,\omega)\), whose bounded-Borel Blackwell game has a value \(u_f(q)\). Define \(c_f=\operatorname{cav}u_f\) by the finite-splitting formula in the canonical statement. The full target is to characterize the payoffs satisfying \(L_f(p)=c_f(p)\) for \textbf{every} \(p\in\Delta(K)\).

Choose \(a<b\) with \(a\le f\le b\), and put \(F=(f-a)/(b-a)\), so \(0\le F\le1\). Write \(F_k(\omega)=F(k,\omega)\) and \(F_q=\sum_kq_kF_k\). The target is unchanged by this normalization. Let \(\Delta_{\mathbb Q}(K)\) denote priors with rational coordinates. All vectors and arrays below are real-valued; no rationality restriction is imposed on witnesses.

Indeed every strategy pair's payoff transforms as \(a+(b-a)\mathbb E F\); positive affine transformations commute with the prescribed supremum and infimum, so \(L_f=a+(b-a)L_F\) and \(u_f=a+(b-a)u_F\). Every splitting has weights summing to one, which also gives \(c_f=a+(b-a)c_F\). This proves the normalization assertion, even for constant \(f\), for which any enclosing bounds \(a<b\) may be used.

\clearpage
\section{The local-geometry certificate condition}

A \textbf{likelihood flow at \(p\)} is an array \(z_k(h)\ge0\) satisfying

\begin{equation}
z_k(\varnothing)=p_k,\qquad
z_k(hij)\text{ is independent of }j,\qquad
\sum_{i\in I}z_k(hij)=z_k(h)
\label{eq:flow}
\end{equation}

for every \(k,h,j\). In particular every entry is at most \(p_k\); if \(z_k(h)=0\), all its children are zero. These are direct algebraic identities, including on histories assigned zero flow.

For a prior \(p\) and integer \(n\ge1\), the condition requires the following quantifier order.

\textbf{First choose one lower-capital package} consisting of an integer \(1\le R\le m\), positive weights \(\lambda_r\), vectors \(q_r\in\Delta(K)\), and arrays \(C_r:H\to[0,1]\), such that

\begin{equation}
\sum_r\lambda_r=1,\qquad \sum_r\lambda_rq_r=p. \label{eq:split}
\end{equation}

At each \(h\) and for each \(r\), form row increment vectors

\[
D^r_i(h)=\big(C_r(hij)-C_r(h)\big)_{j\in J}\in\mathbb R^J.
\]

Require the finite convex-hull intersection

\begin{equation}
\operatorname{conv}\{D^r_i(h):i\in I\}
 \cap\mathbb R_-^J\ne\varnothing \label{eq:lower_hull}
\end{equation}

and the boundary inequality on \textbf{every} infinite path

\begin{equation}
\liminf_{t\to\infty}C_r(\omega|t)\ge1-F_{q_r}(\omega).
\label{eq:lower_boundary}
\end{equation}

Its root reference is \(\beta=\sum_r\lambda_r(1-C_r(\varnothing))\).

\textbf{Then, for every likelihood flow \(z\) satisfying \eqref{eq:flow}, choose upper-capital arrays} \(A_k:H\to[0,1]\). At each \(h\), form column increment vectors

\[
V^z_j(h)=\left(
 \sum_{i\in I}z_k(hij)\,[A_k(hij)-A_k(h)]
 \right)_{k\in K}\in\mathbb R^K.
\]

Require

\begin{equation}
\operatorname{conv}\{V^z_j(h):j\in J\}
 \cap\mathbb R_-^K\ne\varnothing \label{eq:upper_hull}
\end{equation}

and, again on \textbf{every} path and for \textbf{every} state,

\begin{equation}
\liminf_{t\to\infty}A_k(\omega|t)\ge F_k(\omega).
\label{eq:upper_boundary}
\end{equation}

Finally require the root inequality

\begin{equation}
\sum_kp_kA_k(\varnothing)\le\beta+1/n. \label{eq:root_budget}
\end{equation}

The payoff condition is:

\begin{quote}
\textbf{(G)} For each \(p\in\Delta_{\mathbb Q}(K)\) and \(n\ge1\), choose one
package satisfying \eqref{eq:split}--\eqref{eq:lower_boundary}. For every flow
\(z\) satisfying \eqref{eq:flow}, there must then exist capitals \(A_k\)
satisfying \eqref{eq:upper_hull}--\eqref{eq:root_budget}.
\end{quote}

The lower package is fixed before \(z\), whereas the upper arrays may depend
on \(z\). Likelihood flows preserve the informed-strategy quantification.
The criterion uses no explicit behavioral kernels, expected-payoff
optimization, or game values: its inputs are path values of \(F\), and its
tests are conserved flows, finite convex geometry, and pointwise boundary
and root inequalities.

\begin{theorem}[Complete intrinsic criterion]\label{thm:main}
For the fixed finite information and action structure and every bounded Borel \(f\),

\[
\forall p\in\Delta(K),\quad L_f(p)=(\operatorname{cav}u_f)(p)
\quad\Longleftrightarrow\quad \text{(G)}.
\]

The result covers singleton alphabets, simplex faces, arbitrary finite-prefix dependence, and absence of optimal-strategy attainment. The choice of normalization bounds does not affect the equivalence.
\end{theorem}

\section{Established results and the bounded-capital reduction}

Blackwell Borel determinacy \cite[Theorem 10, p.~1578]{Martin} supplies existence of the nonrevealing value: the game has finite action sets, simultaneous moves, complete observation of past joint actions, and bounded Borel payoff \(F_q\). We use it only for these complete-information games. We do not invoke determinacy for the incomplete-information game.

The second established input is a finite-tree Borel superhedging theorem. Assign each \(h\in H\) a nonempty closed convex set \(P_h\subseteq\Delta(X)\). For any Borel \(W:\Omega\to[0,1]\), let

\[
c=\sup_{\pi_h\in P_h\ \forall h}\mathbb E_\pi W,
\]

where the expectation uses the precise probability tree selected by all its local kernels. Then for each \(\delta>0\) there is a globally bounded-below, real-valued array \(M:H\to\mathbb R\) satisfying

\[
\sup_{\pi\in P_h}\sum_{x\in X}\pi_x M(hx)\le M(h),\qquad
\liminf_tM(\omega|t)\ge W(\omega)\quad(\omega\in\Omega),
\]

and \(M(\varnothing)<c+\delta\). This is the combination of definition (3) on p.~314, Theorem 2 on p.~315, and Theorem 13 on p.~318 in T'Joens and De Bock \cite{TDB}. Theorem 2 identifies the axiomatic upper expectation with the game-theoretic superhedging expectation; Theorem 13 identifies the latter with the Borel measure upper expectation. The product sigma field is the Borel sigma field because the alphabet is finite. The theorem has no regularity assumption on the nodewise credal sets beyond the stated finite-dimensional hypotheses.

The capital may always be taken in \([0,1]\). First \(M(h)\ge0\) at every \(h\): otherwise the local inequality and nonemptiness of \(P_h\) give a child with \(M(hx)\le M(h)<0\). Repeating this produces an infinite path whose liminf is negative, contrary to the terminal inequality. Next put \(M'=\min(M,1)\). For every \(\pi\in P_h\), concavity and monotonicity of \(s\mapsto\min(s,1)\) give

\[
\sum_x\pi_x\min(M(hx),1)
\le\min\!\left(\sum_x\pi_xM(hx),1\right)
\le\min(M(h),1).
\]

Its liminf is \(\min(\liminf M,1)\), which still dominates \(W\), and its root cost can only decrease. This proves the claimed bounded-capital version, including terminal domination on paths of probability zero.

\section{Flows and behavioral likelihoods}

Given any informed behavioral strategy \(\sigma\), define

\begin{equation}
z_k(h)=p_k\prod_{s<t}\sigma_k(i_s\mid h_s),
\qquad h=((i_0,j_0),\ldots,(i_{t-1},j_{t-1})).
\label{eq:flow_likelihood}
\end{equation}

These entries satisfy \eqref{eq:flow}. The opponent's likelihood is deliberately absent.

Conversely, given a flow \(z\), choose a fixed distribution \(x^0\in\Delta(I)\) and set

\begin{equation}
\sigma_k(i\mid h)=
\begin{cases}
z_k(hij)/z_k(h),&z_k(h)>0,\\
x^0_i,&z_k(h)=0.
\end{cases}
\label{eq:flow_reconstruction}
\end{equation}

The quotient is independent of \(j\), and the conservation identity makes it a probability vector. Induction proves \eqref{eq:flow_likelihood} at every history. In the zero case all children are zero, so the induction also holds there. At zero-prior states it holds because all entries are zero.

Against any uninformed strategy \(\tau\), the probability of the cylinder fixing state \(k\) and public history \(h\) is exactly

\begin{equation}
z_k(h)\prod_{s<t}\tau_{j_s}(h_s). \label{eq:cylinder_law}
\end{equation}

Thus flows encode the exact statewise likelihood structure with the opponent's current information restriction. The identities are valid off the support as well as on it.

If \(z\) was encoded from a strategy, reconstruction can change that strategy only at state-history nodes with \(z_k(h)=0\). Formula \eqref{eq:cylinder_law} then gives identical state-public cylinder probabilities for the original and reconstructed strategies against every opponent. They therefore induce identical Borel laws and payoff expectations. This also justifies using the reconstruction in upper bounds for an arbitrary originally supplied strategy.

\section{Soundness of each lower package}

For fixed \(r,h\), \eqref{eq:lower_hull} supplies coefficients \(x_i(h)\ge0\), summing to one, with

\begin{equation}
\sum_i x_i(h)C_r(hij)\le C_r(h)\qquad(j\in J). \label{eq:lower_step}
\end{equation}

Choose such coefficients at every history. Since \(H\) is countable, this defines a behavioral nonrevealing strategy without any measurable-selection issue. Against every opponent \(\tau\), the process \(C_r(\omega|t)\) is a bounded nonnegative supermartingale: \eqref{eq:lower_step} holds for each current pure column, and averaging over \(\tau(h)\) preserves it. By \eqref{eq:lower_boundary} and Fatou's lemma,

\[
\mathbb E_{x,\tau}(1-F_{q_r})
\le\mathbb E_{x,\tau}\liminf_t C_r(\omega|t)
\le\liminf_t\mathbb E_{x,\tau}C_r(\omega|t)
\le C_r(\varnothing).
\]

Therefore \(u_F(q_r)\ge1-C_r(\varnothing)\). Weighting the inequalities and using \eqref{eq:split} gives

\begin{equation}
\beta\le\sum_r\lambda_r u_F(q_r)\le c_F(p). \label{eq:reference_bound}
\end{equation}

This bound uses only the package already fixed; it is independent of every subsequently supplied flow.

In fact, for every bounded Borel \(F\) and every prior \(p\), the supremum of the root references \(\beta\) over all packages \eqref{eq:split}, \eqref{eq:lower_hull}, and \eqref{eq:lower_boundary} with \(R\le m\) equals \(c_F(p)\). The displayed bound proves one direction. Section 9 gives an exact maximizing split, and the lower-capital construction in the first half of Section 10 uses neither \(L_F=c_F\) nor any particular flow; applied to that split, it gives \(\beta>c_F(p)-2\delta\) for every \(\delta>0\). This proves the reverse direction. Thus the lower package is itself a value-free geometric representation of concavification, not an unspecified numerical reference supplied by the game values.

\section{Soundness of the upper-capital response}

Fix a flow \(z\), its reconstructed strategy \(\sigma\), and arrays \(A_k\) satisfying \eqref{eq:upper_hull} and \eqref{eq:upper_boundary}. The hull condition supplies a single column distribution \(y(h)\in\Delta(J)\) at every \(h\), with

\begin{equation}
\sum_jy_j(h)\sum_i z_k(hij)
 [A_k(hij)-A_k(h)]\le0\qquad(k\in K). \label{eq:response_weighted_step}
\end{equation}

For \(z_k(h)>0\), division by \(z_k(h)\), followed by \eqref{eq:flow_reconstruction}, gives

\begin{equation}
\sum_{i,j}\sigma_k(i\mid h)y_j(h)A_k(hij)\le A_k(h).
\label{eq:response_step}
\end{equation}

If \(p_k>0\), every positive-probability history under the conditional law in state \(k\) has \(z_k(h)>0\), by \eqref{eq:cylinder_law}. Consequently \(A_k(\omega|t)\) is a bounded nonnegative supermartingale under that conditional law. No condition after division is needed at a state-history cylinder of probability zero. The boundary \eqref{eq:upper_boundary}, which holds everywhere, and Fatou yield

\[
\mathbb E_{k,\sigma_k,y}F_k\le A_k(\varnothing).
\]

Multiply by \(p_k\) and sum over its positive coordinates; zero-prior states contribute nothing. This proves

\begin{equation}
\inf_\tau\mathbb E_{p,\sigma,\tau}F
\le\mathbb E_{p,\sigma,y}F
\le\sum_kp_k A_k(\varnothing). \label{eq:response_bound}
\end{equation}

The same \(y(h)\) works for all states at that node. Its current randomization cannot depend on the hidden state or the current informed action. This information restriction is precisely the common convex combination in \eqref{eq:upper_hull}.

\section{A behavioral proof of the universal lower bound}

For every bounded Borel payoff, \(L_F\ge c_F\). Here is a complete proof that does not rely on tail measurability or public revelation of a split.

Take any finite splitting \(p=\sum_r\lambda_rq_r\), any informed strategies \(\sigma^r\), and a state \(k\) with \(p_k>0\). Put \(\alpha_r(k)=\lambda_rq_{r,k}/p_k\),

\[
L_r(k,h)=\prod_{s<t}\sigma^r_k(i_s\mid h_s),\qquad
D(k,h)=\sum_r\alpha_r(k)L_r(k,h).
\]

At positive \(D(k,h)\) prescribe

\[
\sigma_k(i\mid h)=
\frac{\sum_r\alpha_r(k)L_r(k,h)\sigma^r_k(i\mid h)}{D(k,h)};
\]

at zero \(D\), or zero \(p_k\), prescribe any fixed distribution. The numerator is the next child \(D(k,hi j)\), independent of the current \(j\). Multiplication telescopes, giving \(\prod_{s<t}\sigma_k(i_s\mid h_s)=D(k,h)\) at every history, including zero likelihoods. For any common opponent \(\tau\), its factor \(\prod_{s<t}\tau_{j_s}(h_s)\) is independent of \(r\). Hence the constructed law satisfies, on every state-public cylinder and thus on the generated Borel sigma field,

\begin{equation}
\mathbb P_{p,\sigma,\tau}
 =\sum_r\lambda_r\mathbb P_{q_r,\sigma^r,\tau}. \label{eq:mixture_law}
\end{equation}

For each \(r\), choose \(\sigma^r\) guaranteeing \(L_F(q_r)-\varepsilon\). Such approximate strategies exist by the definition of supremum; no attainment is assumed. Equation \eqref{eq:mixture_law} implies \(L_F(p)\ge\sum_r\lambda_rL_F(q_r)-\varepsilon\). Letting \(\varepsilon\downarrow0\) proves concavity of \(L_F\). Ignoring \(k\) gives \(L_F(q)\ge u_F(q)\) for each \(q\). Therefore \(L_F(p)\ge\sum_r\lambda_ru_F(q_r)\) for every split, and taking their supremum gives

\begin{equation}
L_F(p)\ge c_F(p). \label{eq:universal_lower}
\end{equation}

This is the universal lower bound recorded in Part 1 of the proof of the source theorem \cite{KLS}. The preceding likelihood argument verifies it without granting any extra private-memory convention beyond the stated behavioral strategy space.

\section{Sufficiency at rational priors}

Assume (G), and fix rational \(p\) and integer \(n\ge1\). Take the common package. Every informed \(\sigma\) has a flow \(z\), so (G) supplies upper arrays for it. Equations \eqref{eq:response_bound}, \eqref{eq:root_budget}, and \eqref{eq:reference_bound} give

\[
\inf_\tau\mathbb E_{p,\sigma,\tau}F
\le\sum_kp_kA_k(\varnothing)
\le\beta+1/n\le c_F(p)+1/n.
\]

Taking the supremum over \(\sigma\), and then letting \(n\to\infty\), proves \(L_F(p)\le c_F(p)\). Combine with \eqref{eq:universal_lower} to obtain equality at every rational prior.

\section{A maximizing splitting exists with at most \texorpdfstring{\(|K|\)}{|K|} terms}

For every finite splitting and numbers \(a_r\in\mathbb R\), another splitting using a subset of the same \(q_r\) has at most \(m\) positive weights and has weighted \(a_r\)-sum at least as large. Remove zero weights. If more than \(m\) remain, the \(q_r\in\mathbb R^m\) are linearly dependent. Choose nonzero \(v=(v_r)\) with \(\sum_r v_rq_r=0\). Because every \(q_r\) sums to one, \(\sum_r v_r=0\); therefore \(v\) has both signs. Reverse its sign if necessary so \(\sum_rv_ra_r\ge0\), and set

\[
t=\min_{v_r<0}\frac{\lambda_r}{-v_r}>0,\qquad
\lambda'_r=\lambda_r+tv_r.
\]

The new weights are nonnegative, conserve the barycenter and their sum, make at least one formerly positive weight zero, and do not reduce the objective. Iterating finitely many times leaves at most \(m\) weights. Apply this with \(a_r=u_F(q_r)\). Thus the supremum defining \(c_F(p)\) is unchanged if restricted to at most \(m\) terms.

The function \(u_F\) is 1-Lipschitz and in particular continuous; Section 11 proves this directly without using splitting attainment. Pad any splitting of at most \(m\) terms with zero weights to exactly \(m\) terms. The feasible set

\[
\{(\lambda,q_1,\ldots,q_m)\in\Delta(\{1,\ldots,m\})\times\Delta(K)^m:
\sum_{r=1}^m\lambda_rq_r=p\}
\]

is nonempty and compact, since the barycenter equality is closed. The continuous objective \(\sum_r\lambda_ru_F(q_r)\) attains a maximum there. By support reduction its maximum equals \(c_F(p)\). Discarding the zero weights gives an exact maximizing split with \(1\le R\le m\) positive terms. This finite-dimensional attainment does not assert that any infinite-game strategy attains its optimization.

\section{Necessity: one common lower package, then all flows}

Assume \(L_F(p)=c_F(p)\) for every prior, and fix a prior \(p\) and integer \(n\ge1\). Write \(c=c_F(p)=L_F(p)\) and choose \(\delta=1/(5n)\).

Choose an exact maximizing split with \(R\le m\) and \(\sum_r\lambda_ru_F(q_r)=c\), using Section 9. For each \(r\), choose a nonrevealing strategy \(x^r(h)\in\Delta(I)\) such that

\begin{equation}
\inf_\tau\mathbb E_{x^r,\tau}F_{q_r}\ge u_F(q_r)-\delta.
\label{eq:nonrevealing_guarantee}
\end{equation}

At node \(h\), use the credal set

\[
P^r_h=\{(x^r_i(h)b_j)_{ij}:b\in\Delta(J)\}.
\]

It is nonempty, closed, and convex, as the continuous linear image of a finite simplex. A compatible precise tree corresponds exactly to an uninformed opponent against \(x^r\): its \(J\)-marginal recovers \(b\), and its kernel is the displayed independent product. This statement holds on all nodes; assignments on root-null nodes do not affect the root law. For \(W=1-F_{q_r}\), the robust expectation is

\[
\sup_\tau\mathbb E_{x^r,\tau}(1-F_{q_r})
=1-\inf_\tau\mathbb E_{x^r,\tau}F_{q_r}
\le1-u_F(q_r)+\delta.
\]

The bounded-capital duality of Section 3 supplies \(C_r\in[0,1]\) with \eqref{eq:lower_boundary},

\[
\sum_i x^r_i(h)C_r(hij)\le C_r(h)\quad(j\in J),
\qquad C_r(\varnothing)<1-u_F(q_r)+2\delta.
\]

The coefficients \(x^r(h)\) verify \eqref{eq:lower_hull}. This package depends only on \(p,n\), and its root reference satisfies

\begin{equation}
\beta=\sum_r\lambda_r(1-C_r(\varnothing))
>\sum_r\lambda_r u_F(q_r)-2\delta=c-2\delta.
\label{eq:reference_cost}
\end{equation}

Now let \(z\) be \textbf{any} likelihood flow. Reconstruct its informed strategy \(\sigma\) by \eqref{eq:flow_reconstruction}. Since the infimum against this particular \(\sigma\) is at most the maxmin \(L_F(p)=c\), choose a response \(\tau\) with

\begin{equation}
\gamma:=\mathbb E_{p,\sigma,\tau}F<c+\delta. \label{eq:tailored_response}
\end{equation}

For each state \(k\), apply the duality to the \textbf{singleton} precise credal sets

\[
P^{k}_h=\{(\sigma_k(i\mid h)\tau_j(h))_{ij}\}.
\]

Singletons are nonempty, closed, and convex. Their unique compatible precise tree is the fixed-state-\(k\) kernel law against \(\sigma_k,\tau\), defined independently of whether \(p_k>0\). For \(W=F_k\), obtain \(A_k\in[0,1]\) satisfying \eqref{eq:upper_boundary},

\begin{equation}
\sum_{i,j}\sigma_k(i\mid h)\tau_j(h)A_k(hij)\le A_k(h),
\qquad
A_k(\varnothing)<\mathbb E_{k,\sigma_k,\tau}F_k+\delta.
\label{eq:precise_hedge}
\end{equation}

Multiplying the local inequality by \(z_k(h)\), and using \(z_k(hij)=z_k(h)\sigma_k(i\mid h)\) everywhere, shows that the single common coefficients \(\tau_j(h)\) verify \eqref{eq:upper_hull}. When \(z_k(h)=0\), both sides of this multiplied inequality are zero; no division has occurred. Weighting the root estimates gives

\[
\sum_kp_kA_k(\varnothing)<\gamma+\delta<c+2\delta
<\beta+4\delta<\beta+1/n.
\]

Thus \eqref{eq:root_budget} holds for this flow. The package was already fixed and the argument applies to every flow. Necessity holds at every prior, in particular at all rational priors required in (G).

\section{Extension from rational priors to every prior}

Both \(L_F\) and \(u_F\) are 1-Lipschitz for the \(\ell^1\) distance. For a fixed admissible strategy pair, its expected payoff at \(p\) has the form \(\sum_kp_ka_k\), where \(a_k\in[0,1]\) is the conditional expectation under that pair. It follows that its change is at most \(\|p-p'\|_1\). The strategy spaces are independent of the prior, including prescriptions at zero-prior states. Taking either relevant supremum/infimum preserves the same uniform Lipschitz bound.

Here is a proof that \(c_F=\operatorname{cav}u_F\) is also 1-Lipschitz, including the boundary. Given \(p,p'\), choose a coupling matrix \(T_{ab}\ge0\) with row sums \(p_a\), column sums \(p'_b\), and diagonal entries \(T_{aa}=\min(p_a,p'_a)\). Explicitly, put these entries on the diagonal; the nonnegative residual row and column masses have disjoint supports and equal total \(d=\tfrac12\|p-p'\|_1\). If \(d>0\), couple them by their product divided by \(d\); if \(d=0\), no residual is needed. For \(p_a>0\), let \(S_{ab}=T_{ab}/p_a\), and choose any probability row when \(p_a=0\).

Transport any finite splitting \(p=\sum_r\lambda_rq_r\) by setting \(q'_r=q_r S\). Then \(\sum_r\lambda_rq'_r=p'\). Moreover,

\[
\begin{aligned}
\sum_r\lambda_r\|q'_r-q_r\|_1
&\le\sum_r\lambda_r\sum_aq_{r,a}\|S_{a\cdot}-e_a\|_1\\
&=\sum_ap_a\,2(1-S_{aa})
=2\sum_a(p_a-T_{aa})=\|p-p'\|_1.
\end{aligned}
\]

Terms with \(p_a=0\) vanish: for every positive \(\lambda_r\), barycenter conservation implies \(q_{r,a}=0\). By the Lipschitz property of \(u_F\),

\[
\sum_r\lambda_ru_F(q'_r)
\ge\sum_r\lambda_ru_F(q_r)-\|p-p'\|_1.
\]

Take the supremum over all splits on the right. This gives \(c_F(p')\ge c_F(p)-\|p-p'\|_1\). Interchanging \(p,p'\) proves the opposite bound. Rational priors are dense in the whole simplex. Equality of these two continuous functions at rational priors therefore extends to every prior, completing sufficiency and the theorem.

\section{Scope and semantic audit}

All quantified functionals are bounded Borel; no tail, continuity, prefix, or finite-horizon restriction is added. The terminal constraints hold on every path, including paths assigned probability zero by every relevant root law. The Borel duality theorem is used to prove their \textbf{necessity}, rather than replacing them with almost-sure conditional-expectation statements. Every approximation error is explicit. No optimizing game strategy needs to exist; the finite-dimensional prior splitting does attain its maximum, and this is proved in Section 9.

The condition is structural in the certificate sense: it tests the functional against real arrays by flow conservation, finite convex geometry, and pointwise boundary comparisons. The common column hull encodes the uninformed player's information constraint. The likelihood flows are equivalent to the informed player's likelihoods, as proved algebraically; that interpretation is openly stated, not concealed. The extra capital arrays are dual certificates with local inequalities and global boundary domination, rather than a mere relabeling of the target equality. The result supplies necessity as well as sufficiency and extracts actual response guarantees from any witnesses.

The characterization is infinitary. An arbitrary Borel functional need not have a finite input encoding, and its all-path boundary comparisons are not asserted decidable or finitely checkable. The theorem does not claim a simpler qualitative classification or an effective algorithm. These are scope statements, not additional hypotheses. Finite recognition of an arbitrary Borel functional and simpler qualitative subclass criteria are separate questions.

The condition concerns the maxmin alone. Since the upper arrays may depend on the tested flow, the response distribution extracted from \eqref{eq:upper_hull} may also depend on that flow. There is no common response against all informed strategies. In particular the proof makes no inference that the uninformed minimax equals the maxmin. Known tail-payoff examples where the maxmin equals concavification but a game value fails to exist are compatible with this quantifier order.

\section{Explicit checks against boundary and first-stage examples}

If \(F_k(\omega)=a_k\in[0,1]\) is constant in the public path, choose one split with \(q_1=p\), \(\lambda_1=1\), constant \(C_1=1-\sum_kp_ka_k\), and, for every flow, constant \(A_k=a_k\). All row and column increment vectors vanish. Both boundary conditions hold with equality, and the root inequality holds with zero error. Thus the full certificate includes all statewise constant payoffs, including simplex-face priors and singleton action sets.

The condition does not incorrectly accept every bounded continuous payoff. For \(K=I=J=\{0,1\}\), let \(F\) depend only on the first joint action and have state matrices

\[
A^0=\begin{pmatrix}1&0\\0&0\end{pmatrix},
\qquad A^1=\begin{pmatrix}0&0\\0&1\end{pmatrix}.
\]

Writing \(p=\Pr(k=0)\), the nonrevealing matrix is \(\operatorname{diag}(p,1-p)\). A row-zero probability \(x\) gives column payoffs \(px\) and \((1-p)(1-x)\). Their minimum is maximized at \(x=1-p\) for \(0<p<1\), giving \(u_F(p)=p(1-p)\). The column distribution that chooses column zero with probability \(1-p\) gives both rows this same payoff, proving the matching upper bound. At endpoints the value is zero. The function \(p(1-p)\) is concave, so its concavification is itself.

The informed choice of row zero in state zero and row one in state one gives payoffs \(p\) and \(1-p\) against the two columns. Conversely pure column zero bounds every informed strategy by \(p\), and pure column one bounds it by \(1-p\). Therefore \(L_F(p)=\min(p,1-p)\). At \(p=1/2\), the type-dependent first action gives expected payoff \(1/2\) against every opponent. Its associated flow, tested in (G), forces every upper root cost to be at least \(1/2\), by Section 6. Every lower-package reference is at most \(1/4\), by Section 5. Hence \eqref{eq:root_budget} is impossible for \(n\ge5\). This is an exact algebraic exclusion, without a finite computational substitute for any infinite-game assertion.

\begin{thebibliography}{9}
\bibitem{KLS}
Gil Bar Castellon Koltun, Ehud Lehrer, and Eilon Solan.
\emph{The Maxmin Value of Repeated Games with Incomplete Information on
One Side and Tail-Measurable Payoffs}. arXiv:2512.01581 (2025).
Definition 2.1, p.~6, specifies the game; Remark 2.2 and Definitions 2.3--2.4,
p.~7, specify the evaluation and nonrevealing game. Theorem 2.5, p.~8,
states the tail-payoff theorem; Section 3, Part 1 of its proof, p.~8, records
the universal lower bound. Section 5, p.~18, asks the characterization question.
\url{https://arxiv.org/pdf/2512.01581v1}.

\bibitem{Martin}
Donald A. Martin. \emph{The Determinacy of Blackwell Games}.
Journal of Symbolic Logic \textbf{63}(4), 1565--1581 (1998).
The behavioral finite-action setup is given on p.~1567. Theorem 10,
p.~1578, supplies the bounded-Borel Blackwell determinacy used only for the
complete-information nonrevealing games.
\url{https://doi.org/10.2307/2586667}.
\url{https://www1.cmc.edu/pages/faculty/MONeill/math188/papers/martin2.pdf}.

\bibitem{TDB}
Natan T'Joens and Jasper De Bock.
\emph{Global Upper Expectations for Discrete-Time Stochastic Processes:
In Practice, They Are All The Same!}
Proceedings of Machine Learning Research \textbf{147}, 310--319 (2021).
Definition (3), p.~314; Theorem 2, p.~315; Theorem 13, p.~318.
\url{https://proceedings.mlr.press/v147/t-joens21a/t-joens21a.pdf}.
\end{thebibliography}
\end{document}
